\scanpage{39}
\mspage{19}
and $j:U-Y$ induced by $\tau_u$, I say that $H_1(i)=H_1(j)$, or what amounts to the same, since $Y\xrightarrow{\alpha}J_Y^1$ induces an isomorphism $H_1(\alpha):H_1(Y)-H_1(J_Y^1)$, that $\unsure{H_1(\alpha)|H'_1(\alpha)|H_1^*(\alpha)}\unsure{H_1(i)|H'_1(i)}=H_1(\alpha)H_1(j)$, or again $H_1(\alpha i)=H_1(\alpha j)$.

Now we have $\alpha j=\tau_u(\alpha i)$ therefore \deleted{it suffices to}
\[
H_1(\alpha j)=H_1(\tau_u)H_1(\alpha i)
\]
and we know that $H_1(\tau_u)=\operatorname{id}$, OK. That being so, let $I=\widehat Y\setminus Y$, $J=\widehat Y\setminus U$, in such a way that $J\supset I$, \unsure{and|\deleted{and}} Let $f:\widehat Y\simeq\widehat Y$ be obtained from $j:U\subset Y$ by passing to the completions, therefore \inserted{$f(U)=\widehat Y-f(J)\subset Y=\widehat Y-I$} \deleted{\unsure{we have} $\unsure{f(J))|f(J)}$} implies $f(J)\supset I$ i.e. $J\supset f^{-1}(I)$.

We see that the passage from $H_1(U)$ to $\unsure{H^1(Y)|H_1(Y)|H'(Y)}$, by filling in holes at $J\setminus I$, has as kernel the \inserted{sum of the} $L_j$ with $j\in J\setminus I$, and for the analogous reason the kernel of $H_1(U)\unsure{-|\to}H_1(Y)$ induced by $j$ is the sum of the $L_j$ with $j\in J\setminus f(I)$, moreover the $L_j$ $(j\in J)$ in $H_1(U)$ are $\simeq\widehat{\mathbb Z}$ (resp. $\widehat{\mathbb Z}'$ in characteristic $p>0$\unsure{..|...}) and the sub-$\widehat{\mathbb Z}$-module generated is $\widehat{\mathbb Z}^{J}/\widehat{\mathbb Z}$ (diagonal), therefore \deleted{if $\operatorname{card}J\geq3$\unsure{,}} the submodules generated by two subsets \inserted{$I,I'$} of $J$ distinct from $J$, can be equal only if $I=I'$, or $I$ and $I'$ of cardinality $\mu-1$ (where $\mu=\operatorname{card}J$)
\inserted{But here after possibly shrinking $U$, we may assume}
\margin{38}
